% Mathematical contribution and independent audit for the parent report.
% This is an intermediate LaTeX fragment, not a standalone deliverable.
% Environments theorem, lemma, proposition, corollary, remark are assumed.
% The named predecessor is the incoming witness_discharge package at the
% parent's pinned repository snapshot.  Its source should be cited by parent.

\section{Deterministic equality at witness discharge}
\label{dd:main}

The equality-aware witness-discharge presentation has three instruction
families, but its discharge instruction still asks for an equality proof and
a substitution target.  Both can be removed.  The opened clause itself can
determine one equality to use; all other instances use the identity
substitution.  The resulting instruction has only its major proof and one
branch proof.  Its target is computed from the major formula.

This change weakens a primitive inference.  It neither adds an arithmetic
axiom nor changes the meaning of the paired binder.  In particular, the
substitution target is not an unchecked hint, a hidden proof argument, or a
search result: the following fixed syntactic prescription determines it.

\subsection{Canonical syntax and the complete rule}

Retain the full computation-normal form syntax of the predecessor.  Terms
are the normal forms of
\[
 t ::= x\mid 0\mid S(t)\mid I(n,a;u.h)
\]
under the two context-closed reductions
\[
 I(0,a;u.h)\longrightarrow a,
 \qquad
 I(Sn,a;u.h)\longrightarrow h[I(n,a;u.h)/u].
\]
Only $u$ is bound in $h$.  Write $N$ for normalization.  Every substitution
is capture avoiding and is followed by $N$ on numerical subterms.  Formulas
are equations and $D_{x,y}B$; normalization does not remove formula binders.
Set $H=(0=0)$ and $1=S0$, $2=S(S0)$.

A clause $[\vec A;y.B(y)]$ has zero or one antecedent, and $y$ is absent
from the antecedent.  Its ordinary reading is $\forall y\,B(y)$ when
$\vec A$ is empty and $A\to\forall y\,B(y)$ otherwise.  An ordinary
assumption $F$ is $[\,;v.F]$, where $v$ is unused.  As before, the initial
context contains the one fixed axiom clause
\[
             \mathcal R=[\,;y.y=y].
\]
Thus reflexivity is an explicitly counted axiom, available by clause use.
The ordinary reading of $D$ remains
\[
 (D_{x,y}B)^{\circ}
   =\exists x\bigl(B^{\circ}(x,0)\to
                       \forall y\,B^{\circ}(x,y)\bigr).
\]
The designated contradiction is still $0=S0$.

For a fresh parameter $b$ and a matrix $B(x,y)$, compute
\[
 B_0=N(B(b,0)),\qquad B_1=N(B(b,1)),
\]
and define the term
\[
 \rho(B,b)=
 \begin{cases}
 t,& B_0\equiv(0=0)\ \text{and}\ B_1\equiv(b=t),\\
 b,& \text{otherwise}.
 \end{cases}
 \tag{\(\rho\)}\label{dd:rho}
\]
The tests are literal comparisons of canonical syntax.  In particular, the
successful case requires the left side of $B_1$ to be exactly the newly
opened parameter.  No equation recognition up to provability is intended.
The extracted right side $t$ may itself contain $b$.

The new presentation $\mathsf{DD}$ has the following three instruction
families.
\begin{enumerate}
\item \textbf{Use.} For $[\vec A;y.B]\in\Gamma$, prove the single
antecedent when it is present and infer $N(B(t))$ at any well-scoped $t$.
An empty-antecedent clause has no proof premise.
\item \textbf{Introduce.} From
\[
 \Gamma,B(t,\eta)\vdash B(t,Sz)
 \quad\text{infer}\quad
 \Gamma\vdash D_{x,y}B(x,y),
\]
where $z$ is fresh and $\eta=z$ precisely when $B$ has equation root and
the canonical witness $t$ is literally $0$; otherwise $\eta=0$.
\item \textbf{Discharge.} Write
\[
  \mathcal B_b=[[N(B(b,0))];y.B(b,y)].
\]
From
\[
 \Gamma\vdash D_{x,y}B(x,y),\qquad
 \Gamma,\mathcal B_b\vdash F
\]
infer
\[
       \Gamma\vdash N(F[\rho(B,b)/b]),
 \tag{DD}\label{dd:rule}
\]
provided the displayed exported formula has no free $b$.
The parameter $b$ is fresh for $\Gamma$ and the passive parameters of $B$.
\end{enumerate}
The substitution in DD is simultaneous: it does not recursively substitute
inside inserted copies of $\rho(B,b)$.  The free-variable test is made after
normalization.  There is no equality-proof premise and no user-supplied
target term in DD.  These are all the primitive rules.  No separate
computation, transport, paired-elimination, cut, weakening, or induction
rule is added.  Computation normalization remains part of the trusted
syntax and must be counted in every complete-checker comparison.

\subsection{Soundness and the ordinary elimination rule}

\begin{lemma}[The determined equality is available in its own branch]
\label{dd:local-equality}
For every opened matrix, $\Gamma,\mathcal B_b$ proves
$b=\rho(B,b)$ by Use alone, with the fixed reflexivity clause $\mathcal R$.
\end{lemma}
\begin{proof}
In the successful case, use $\mathcal R$ at $0$ to prove $H$ and use
$\mathcal B_b$ at $1$.  Its antecedent is literally $H$ and its instance
is literally $b=t$.  In the fallback case, use $\mathcal R$ at $b$.
No equality-transport rule or arithmetic induction is used in either case.
\end{proof}

\begin{theorem}[Soundness of deterministic discharge]
\label{dd:sound}
Rule DD is sound in intuitionistic first-order equality logic over the
ordinary scalar theory, including for an extracted target containing $b$.
It is also a derived inference of the predecessor's equality-aware
witness-discharge presentation.
\end{theorem}
\begin{proof}
Open the ordinary existential of $D_{x,y}B$ at fresh $b$.  Its local
assumption is the reading of $\mathcal B_b$.  Lemma~\ref{dd:local-equality}
supplies $b=\rho(B,b)$, and the given branch supplies $F$.  Equality
substitution gives $F[\rho(B,b)/b]$.  The computation equations identify
this formula with its normal form.  The required absence of $b$ from that
normal form permits ordinary existential elimination.

Equality substitution is valid when its target contains its source
parameter: the substitution is simultaneous.  It is the exported formula,
not the target term by itself, that must be independent of the discharged
parameter.  For the predecessor's rule, insert exactly the equality proof
of Lemma~\ref{dd:local-equality} as its omitted equality branch, and supply
its determined target $\rho(B,b)$.  The end formula is unchanged.
\end{proof}

\begin{lemma}[Ordinary paired elimination]
\label{dd:ordinary-open}
The rule
\[
 \frac{\Gamma\vdash D_{x,y}B\qquad
       \Gamma,\mathcal B_b\vdash A}
      {\Gamma\vdash A},
 \qquad b\notin\operatorname{FV}(\Gamma,A),
\]
is derivable in $\mathsf{DD}$.
\end{lemma}
\begin{proof}
Apply DD to the given branch.  Since $A$ has no free $b$, simultaneous
substitution leaves it unchanged for every value of $\rho(B,b)$.
This remains true when extraction is forced and when the extracted target
contains $b$.  Therefore $N(A[\rho(B,b)/b])\equiv A$.
\end{proof}

This observation is essential: ordinary elimination does not require a
certificate flag that disables extraction.  Independence of the exported
formula already makes the determined equality harmless.

Weakening, injective fresh renaming, and cut of an ordinary assumption are
admissible by direct induction on finite derivations.  For DD the matrix
and hence its determined target do not change when a proof of a retained
ordinary assumption is inserted.  Freshen the local parameters before
inserting proofs.  An occurrence of the unconditional constant clause for
the cut assumption is replaced by its supplied proof.  These structural
arguments require neither general term substitution in derivations nor
equality transport.

\subsection{Numerical cases before any equality transport}

For an unused witness variable $w$, abbreviate
$q_yA(y)=D_{w,y}A(y)$.  The following derived rules are available before
recovering equality:
\[
 \frac{\Gamma,A(0)\vdash A(Sz)}{\Gamma\vdash q_yA(y)},
 \qquad
 \frac{\Gamma\vdash q_yA(y)\qquad\Gamma\vdash A(0)}
      {\Gamma\vdash A(t)}.
 \tag{Cases}\label{dd:cases}
\]
For the first rule use witness $1$ in Introduce; its mode is cases at every
matrix root.  For the second, use Lemma~\ref{dd:ordinary-open}; the opened
clause is $[A(0);y.A(y)]$ because $w$ is unused, and Use gives $A(t)$.
Consequently, proofs of $A(0)$ and $A(Sz)$ with fresh $z$ give a proof of
$A(t)$ for every numerical term $t$.  All needed uses of ordinary paired
elimination remain valid under the forced extractor by the preceding lemma.

\subsection{Unrestricted equality from the fixed probe at one}

Use the inherited normalizing selector and parity terms
\[
 C(v,a,c)=N(I(v,a;u.c)),\qquad
 P(y)=N(I(y,0;u.C(u,1,0))),
\]
where the selector body $c$ does not contain its displayed bound variable
$u$.  Their relevant identities are identities of canonical syntax:
\[
 C(0,a,c)\equiv a,\quad C(Sv,a,c)\equiv c,
 \quad P(1)\equiv1,\quad P(2)\equiv0.
\]
No arithmetic range theorem for $P$ and no induction on $P$ is required.

\begin{theorem}[Uncrossed term transport]
\label{dd:term-transport}
For every numerical term template $h(w)$, the inference
\[
 \frac{\Gamma\vdash s=t\qquad\Gamma\vdash N(h(s))=r}
      {\Gamma\vdash N(h(t))=r}
\]
is derivable in $\mathsf{DD}$.
\end{theorem}
\begin{proof}
Choose $x,y$ fresh for all passive data and set
\[
 B(x,y)=\left(
 C(y,0,C(P(y),N(h(x)),x))
 =C(y,0,C(P(y),r,t))\right).
 \tag{Package}\label{dd:package}
\]
The normalized endpoints are
\[
 B(x,0)\equiv H,\qquad
 B(x,1)\equiv(x=t),\qquad
 B(x,2)\equiv(N(h(x))=r).
\]
Introduce $D_{x,y}B$ at witness $s$.  Its successor goal is
\[
 C(P(Sz),N(h(s)),s)=C(P(Sz),r,t).
\]
Apply the numerical-case construction~\eqref{dd:cases}, at target $P(Sz)$,
to the formula
\[
 A(v)=\bigl(C(v,N(h(s)),s)=C(v,r,t)\bigr).
\]
Its zero instance is the supplied equation $N(h(s))=r$; every successor
instance is the supplied equation $s=t$.  Choose the case eigenparameter
independently of $z$ and fresh for the current retained context.  This proof
ignores Introduce's temporary hypothesis, so it is valid whether $s$ is
literal zero and the mode is progress, or the mode is cases.

Open the resulting $D$ formula at fresh $b$.  Its base is literally $H$,
and its instance at $1$ is literally $b=t$.  Thus the fixed prescription
extracts $\rho(B,b)=t$.  Use the opened clause at $2$, supplied with
reflexivity of $0$, to prove the branch formula
$F=(N(h(b))=r)$.  DD exports
\[
 N\bigl((N(h(b))=r)[t/b]\bigr)
       \equiv (N(h(t))=r).
\]
The last identity is normalization--substitution coherence, rather than
an equality inference inside the theory.  The fresh $b$ does not occur in
$t$, $r$, $h$, or $\Gamma$, so the end formula is outer scoped.  There is
no equality-proof branch at this discharge: the fixed probe at $1$ has
selected its target from the package.
\end{proof}

The order of the two nonzero rows is deliberate.  The predecessor's naming
package can put $b=t$ at $2$; exchanging its two selector branches puts that
equation at the smaller fixed probe $1$.  The formula branch then uses $2$.
The mathematical weakening is independent of this minor convention, but
the complete checker should implement the stated convention exactly.

\begin{corollary}[All ordinary equality reasoning]
\label{dd:full-equality}
The calculus derives symmetry, transitivity, term congruence, crossed
transport, and transport through every Eq/$D$ formula template.
\end{corollary}
\begin{proof}
Take $h(w)=w$ and the reflexive minor $s=s$ in
Theorem~\ref{dd:term-transport}; from $s=t$ it yields $t=s$.  Applying this
symmetry to the theorem's conclusion gives crossed transport.  Congruence
follows by taking the reflexive minor $h(s)=h(s)$ and reversing the
result.  Symmetry followed by transport at the identity template gives
transitivity.  Applying term transport to each side, with symmetry between
the applications, transports an entire equation template.

For a compound template $D_{x,y}A(v,x,y)$ use external induction on formula
depth.  Given $s=t$ and $D_{x,y}A(s,x,y)$, use ordinary paired elimination
at fresh $c$.  Introduce $D_{x,y}A(t,x,y)$ at witness $c$.  This witness is
a parameter rather than literal zero, so Introduce uses cases.  Under its
temporary assumption $A(t,c,0)$, transport backward at the smaller matrix
using $t=s$, obtaining $A(s,c,0)$.  Use the old opened clause at $Sz$ to
obtain $A(s,c,Sz)$, and transport forward at the smaller matrix to
$A(t,c,Sz)$.  Close Introduce and ordinary elimination.  Every recursive
transport acts on a strictly smaller formula.  This proof never applies
induction to a quantified formula.
\end{proof}

\subsection{Exact equivalence with equality-aware discharge}

Let $\mathsf{WD}$ denote the predecessor on the same canonical terms,
formulas, clause contexts, fixed reflexivity axiom, and introduction rule.
Its discharge rule instead supplies the two branch proofs $b=t$ and $F$,
and exports $N(F[t/b])$ if that result is independent of $b$.

\begin{theorem}[Identity on every end sequent]
\label{dd:equivalence}
For every finite canonical context $\Delta$ of zero- or one-antecedent
clauses and every canonical formula $F$,
\[
 \mathcal R,\Delta\vdash_{\mathsf{DD}}F
 \quad\Longleftrightarrow\quad
 \mathcal R,\Delta\vdash_{\mathsf{WD}}F.
\]
Both directions are effective transformations of finite proof trees and
leave the entire end sequent unchanged.
\end{theorem}
\begin{proof}
From DD to WD, translate the given major and branch recursively and insert
the equality proof of Lemma~\ref{dd:local-equality}.  The target of WD is
the computed term $\rho(B,b)$, and the result is literally the same as DD.
Use and Introduce are unchanged.

Conversely, translate a WD derivation recursively.  At a discharge node,
ordinary paired elimination is available by Lemma~\ref{dd:ordinary-open}.
Within its opened context, the recursively translated branches prove
$b=t$ and $F$.  Corollary~\ref{dd:full-equality} gives
$N(F[t/b])$.  The original WD node certifies that this formula is free of
$b$, so ordinary paired elimination exports it unchanged.  The transport
derivation in the corollary was constructed using DD, numerical cases,
and Use, without an untranslated WD node.  Thus this recursive translation
has no circular dependency.  Freshen local parameters at each insertion.
\end{proof}

Composing this identity-language equivalence with the predecessor's
canonical/raw comparison and the native/scalar/iPRA bi-interpretation gives
full first-order intuitionistic primitive recursive arithmetic.  Both
numerical domains remain whole, interpreted equality is actual equality,
and the inherited numerical composite isomorphisms still have graph $x=y$.
Both formula round trips are preserved, because this comparison changes
neither the formulas nor their ordinary readings.  The unchanged
contradiction $0=S0$ and the two effective proof translations also give
equiconsistency.  These conclusions do not rely on standard-model agreement
alone and do not assert that iPRA proves its own consistency.

\subsection{The new rule and numerical substitution}

\begin{proposition}[The extractor introduces no substitution repair]
\label{dd:subst}
Consider a single valid DD inference and a capture-avoiding substitution
$\sigma$ for its passive free parameters, leaving its fresh local parameter
$b$ protected.  After normalized substitution in its premises, the same
DD instruction gives the normalized substitution instance of its original
conclusion.  This statement is about this rule; it does not assert that
the inherited witness-tested introduction keeps its mode under every
substitution.
\end{proposition}
\begin{proof}
If extraction originally succeeds, the normalized base $H$ remains $H$.
The normalized instance $b=t$ becomes $b=N(t\sigma)$, by coherence and
the fact that $b$ is protected.  Extraction therefore still succeeds and
chooses the substituted target.  The two conclusions agree by coherence
of capture-avoiding substitution and normalization.

If extraction originally fails, the original target is $b$, so its
exported canonical formula is the branch formula itself.  The side
condition therefore implies that this branch formula has no free $b$.
Its substitution instance remains independent of $b$.  Whether or not
substitution creates a successful extraction, substituting the newly
determined target for $b$ leaves this branch formula unchanged.  It is
exactly the substitution instance of the old conclusion.
\end{proof}

The successful guard cannot be destroyed by passive-parameter substitution,
but it can be created.  The preceding argument addresses that latter case
explicitly; unconditional preservation of a syntactic test would be false.
Admissibility of numerical substitution for complete derivations still uses
the inherited repair for Introduce when normalization turns a witness into
literal zero, or equivalently the established comparison with the reference
calculus.

\subsection{Proof-checking and the scope of the improvement}

The new discharge checks two proof subtrees, not three, and its certificate
contains no explicit target term.  It adds two fixed normalized instances
of the already supplied matrix, the literal tests in~\eqref{dd:rho}, and
the same normalized substitution and export-scope check as its predecessor.
These changes must all be counted in a matched whole-checker measurement.
The mathematical weakening alone does not guarantee a reduction in every
source-size coordinate or a repository-wide minimum.

Full normalization is a trusted metasyntactic operation, not an
object-language universal evaluator.  No uniform primitive-recursive
normalization bound is asserted.  A fuel-bounded verifier can debit every
normalization and recursive-check operation and reject on exhausted fuel.
Any accepted finite proof has a sufficiently large finite fuel value;
consequently this bounded interface presents the same theorems.  The
additional calls needed to compute both probe instances must be charged to
that same fuel account.  These resource facts neither add an arithmetic
totality axiom nor change either interpretation domain.

The specifically new result is the removal of the explicit equality-proof
premise and target from witness discharge by a fixed, sound probe of the
opened clause.  The canonical iterator syntax, paired binder, clause
interface, fixed reflexivity axiom, two-row naming method, and arithmetic
bi-interpretation are inherited and should be attributed as such.  The
argument here is a written universal proof; finite implementation tests
would not constitute a new proof-assistant verification.
